\documentclass[leqno, 12pt]{jsarticle}
\usepackage{amssymb}
\usepackage{amsthm}
\usepackage{amsmath}
\usepackage{comment}
	%可換図式用
		%\usepackage[all]{xy}
	%重くなるので使わいときはコメント化しておく。
%定理環境 thmはあまり使わずpropをなるべく使う。
%主張は基本的に証明の中で使い、そのため番号を付けない。
%注意環境を付けてみた。
\newtheorem{thm}{定理}
\newtheorem{prop}[thm]{命題}
\newtheorem{cor}[thm]{系}
\newtheorem{lem}[thm]{補題}
\newtheorem{df}[thm]{定義}
\newtheorem{exam}[thm]{例}
\newtheorem{fact}[thm]{事実}
\newtheorem*{claim}{主張}
\newtheorem*{rmk}{注意}
\renewcommand{\proofname}{{\bf 証明}}
%矢印たち。
%\toは\rightarrowと同じ。\getsは\leftarrowと同じ。同値は\iff
%上向き矢はuparrow逆はdownarrow
\newcommand{\To}{\Longrightarrow}
\newcommand{\Gets}{\Longleftarrow}
%黒板太字
\newcommand{\nn}{\mathbb{N}}
\newcommand{\zz}{\mathbb{Z}}
\newcommand{\qq}{\mathbb{Q}}
\newcommand{\rr}{\mathbb{R}}
\newcommand{\cc}{\mathbb{C}}
%無限大
\newcommand{\mgn}{\infty}
%ギリシャン
\newcommand{\dpst}{\displaystyle}
\newcommand{\ep}{\varepsilon}
\newcommand{\al}{\alpha}
\newcommand{\be}{\beta}
\newcommand{\de}{\delta}
\newcommand{\la}{\lambda}
\newcommand{\La}{\Lambda}
\newcommand{\ga}{\gamma}
\newcommand{\lil}{\la\in \La}
%商集合
\newcommand{\qt}[2]{#1/\! #2}
%enumerate環境
\renewcommand{\labelenumi}{{\rm (\arabic{enumi})}}
%以降alignじゃなくてenumerate を使え。
%なんか演算
\DeclareMathOperator{\card}{card}
\DeclareMathOperator{\supp}{supp}
%なんか演算２
\DeclareMathOperator{\bd}{Bdry}
\DeclareMathOperator{\cl}{CL}
\DeclareMathOperator{\nai}{Int}
\DeclareMathOperator{\di}{diam}

%差集合は\setminusで統一する。等号の否定は\neq
%真部分集合は\subsetneqq　偏微分は\partial
%ディスプレイスタイル。\dfracでディスプレイ分数
%空集合は\Oを使う！！！！！！！！！！！！

\title{特別な積分の球値}
\author{ナンブ　キトラ}
\date{\today}
			\begin{document}
\maketitle
この文章では特殊な積分の値を求める．
この文章では積分は全てルベーグ積分である．
\section{ディリクレ積分$\displaystyle \int_0^{\infty}\frac{\sin x}{x}\ dx=\frac{\pi}{2}$}
積分の計算方法として
Feynman's integral trickもしくは
Differentiation under integral sign
と呼ばれるものがある．
これは積分$\int_{X}f(x)\ dx$を求める際に，
$K(x, 0)=1$となる関数$K(x, t)$を用いて
$I(t)=\int_{X}K(x, t)f(x) \ dx$と定義して，
$I$が満たす微分方程式を解いて$I(t)$を求めた上で，$I(0)=\int_{X}f(x)$
を得る方法である．
この方法は物理学者ファインマンの自伝「ご冗談でしょ，ファインマンさん」\cite{2} の「毛色の違った道具」という章で言及され，
よくわからん積分の問題が出たらファインマンのところに
お鉢が回ってきて
この方法で解いたということが記されている．
こうしたことから
ファインマンの名前がついているようだが，\cite{2}の同じ章で言及されている通り
ファインマンはWoodsの解析の教科書 \cite{3}
を参考にこの方法を身につけており，
実際Woodsの本にもこの方法が記されている．

この計算方法の例としてよく取り上げられるのが，ディリクレ積分と呼ばれる
$\int_{0}^{\infty}\sin x/x\ dx$を求める次の方法である．
数学的に怪しい部分があるものをまず最初に紹介する．
\begin{proof}[\textbf{Feynman's integral trickを用いたディリクレ積分の球値？}]
各$t\in [0, \infty)$に対して
\[
F(t)=\int_{0}^{\infty}\exp(-tx)\frac{\sin x}{x}\ dx
\]
と定義する．
このとき
微分と積分を交換して
\[
\frac{d}{dt}F(t)=\int_{0}^{\infty}\frac{\partial }{\partial t}\left(\exp(-tx)\frac{\sin x}{x}\right)\ dx=-\int_{0}^{\infty}\exp(-tx)\sin(x)\ dx
\]
を得る．
このとき
\[
F'(t)=-\int_{0}^{\infty}\exp(-tx)\sin(x)\ dx=\left[\frac{\cos x+t\sin x}{1+t^2}\exp(-tx)\right]_{x=0}^{\infty}=-\frac{1}{1+t^2}
\]
なので，$F(\infty)=0$を用いて
\[
F(t)=F(t)-F(\infty)=\int_{\infty}^tF'(s)\ ds=\int_{\infty}^t\frac{-1}{1+t^2}\ ds
=\frac{\pi}{2}-\arctan t
\]
となる．よって
\begin{align*}
	\int_0^{\infty}\frac{\sin x}{x}\ dx
	&=
	\int_{0}^{\infty}\lim_{t\to 0}\exp(-tx)\frac{\sin x}{x}\ dx
	=\lim_{t\to 0}\int_{0}^{\infty}\exp(-tx)\frac{\sin x}{x}\ dx
	\\
	&=\lim_{t\to 0}F(t)
	=
	\lim_{t\to 0}\left(\frac{\pi}{2}-\arctan t\right)
	=\frac{\pi}{2}
\end{align*}
となる．
\end{proof}

上の議論は計算結果だけ見ると正しいが，
実は数学的に不確かな部分がある．
微分と積分を交換した部分ではない．
その部分はルベーグの収束定理を用いると
簡単に正当化できる．
どこが微妙なのかと言えば
\[
\int_0^{\infty}\frac{\sin x}{x}\ dx
	=
	\int_{0}^{\infty}\lim_{t\to 0}\exp(-tx)\frac{\sin x}{x}\ dx
\]
の部分である．
実は$x$の関数$\sin x/x$はルベーグ可積分ではなく，
$\int_{\rr}|\sin x|/|x|\ dx=\infty$となる．
よってルベーグの収束定理が適用できず，
この積分と極限の記号の交換の正当性は微妙である．
そもそも$\sin x/x$が可積分ではないのに
ルベーグ積分を基礎にした場合，ディリクレ積分は一体どのように定義されているのかという問題もある．ここでは
\[
\int_0^{\infty}\frac{\sin x}{x}\ dx=\frac{\pi}{2}
\]
というのは
\[
\lim_{R\to \infty}\int_0^{R}\frac{\sin x}{x}\ dx=\frac{\pi}{2}
\]
という意味だとすることにして，この値を求める方法を考えよう．
（広義リーマン積分だと思えば別に間違いというわけでもない）．
これを踏まえて改めてFeyman's integral trickを用いて証明しよう．

\begin{proof}[\textbf{Feynman's integral trickを用いたディリクレ積分の球値}]
任意の$R\in [1, \infty)$と$t\in [0, \infty)$に対して
\[
F_R(t)=\int_{0}^{R}\exp(-tx)\frac{\sin x}{x}\ dx
\]
と定義する．
任意の$R$について
$t\in (0, \infty)$のとき
微分と積分を交換して
\[
	\frac{d}{dt}F_R(t)
	=
	\int_{0}^{R}
	\frac{\partial }{\partial t}
	\left(\exp(-tx)\frac{\sin x}{x}\right)\ dx
	=
	-\int_{0}^{R}\exp(-tx)\sin(x)\ dx
\]
を得る．
この等式はルベーグの収束定理を用いると正当化できる．
さて，
\begin{align*}
	\frac{d}{dt}F_R(t)
	&=-\int_{0}^{R}
	\exp(-tx)\sin(x)\ dx
	=
	\left[\frac{\cos x+t\sin x}{1+t^2}\exp(-tx)\right]_{x=0}^{x=R}
	\\
	&=
	\frac{\cos R+t\sin R}{1+t^2}\exp(-tR)
	-\frac{1}{1+t^2}
\end{align*}
なので，$F_R(\infty)=0$を用いて
\begin{align*}
	F_R(t)
	&=
	F_R(t)-F_R(\infty)
	=
	\int_{\infty}^tF_R'(s)\ ds
	=
	\int_{\infty}^{t}\frac{\cos R+s\sin R}{1+s^2}\exp(-sR)\ ds
	+
	\int_{\infty}^t\frac{-1}{1+s^2}\ ds
	\\
	&=
	\int_{\infty}^{t}\frac{\cos R+s\sin R}{1+s^2}\exp(-sR)\ ds
	+
	\frac{\pi}{2}-\arctan t
\end{align*}
となる．
任意の$R$について$[0, R]$上で$\sin x/x$は可積分だから
ルベーグの収束定理から
\[
	\int_0^{R}\frac{\sin x}{x}\ dx
	=
	\lim_{t\to  0}\int_{0}^{R}\exp(-tx)\frac{\sin x}{x}\ dx
\]
がわかる．
よって
\begin{align*}
	\int_0^{R}\frac{\sin x}{x}\ dx
	&=
	\lim_{t\to 0}\int_{0}^{\infty}\exp(-tx)\frac{\sin x}{x}\ dx
	=\lim_{t\to 0}F_R(t)
	\\
	&=
	\lim_{t\to 0}
	\left(
	\int_{\infty}^{t}\frac{\cos R+s\sin R}{1+s^2}\exp(-sR)\ ds
	+
	\frac{\pi}{2}-\arctan t
	\right)
	\\
	&=
	\lim_{t\to 0}
	\left(
	\int_{\infty}^{t}\frac{\cos R+s\sin R}{1+s^2}\exp(-sR)\ ds
	+
	\frac{\pi}{2}
	\right)\\
	&=
	\frac{\pi}{2}+
	\lim_{t\to 0}
	\int_{\infty}^{t}\frac{\cos R+s\sin R}{1+s^2}\exp(-sR)\ ds
\end{align*}
となる．
ここで最右辺の第二項を調べてみよう．
まずルベーグの収束定理から，
\[
	\lim_{t\to 0}
	\int_{\infty}^{t}\frac{\cos R+s\sin R}{1+s^2}\exp(-sR)\ ds
	=
	\int_{\infty}^{0}\frac{\cos R+s\sin R}{1+s^2}\exp(-sR)\ ds
\]
である．よって
\begin{align*}
	\left|
	\int_{\infty}^{0}\frac{\cos R+s\sin R}{1+s^2}\exp(-sR)\ ds
	\right|
	&\le 
	\int_{0}^{\infty}(1+s)\exp(-sR)\ ds
	\\
	&\le 
	\frac{1}{R}+\frac{1}{R^2}
\end{align*}

以上から，
\[
	\left|
	\int_0^{R}\frac{\sin x}{x}\ dx
	-\frac{\pi}{2}
	\right|
	\le 
	\left|
	\int_{\infty}^{0}\frac{\cos R+s\sin R}{1+s^2}\exp(-sR)\ ds
	\right|
	\le 
	\frac{1}{R}+\frac{1}{R^2}
\]
がわかるので，$R\to \infty$とすると
\[
\lim_{R\to \infty}\int_0^{R}\frac{\sin x}{x}\ dx=\frac{\pi}{2}
\]
を得る．これで証明が終わる．
\end{proof}

\section{ガウス積分}
ガウス積分とは
\[
\int_{\rr}\exp(-x^2)\ dx
\]
のことである．有名なこの積分の値は$\sqrt{\pi}$
になる．
大学で習うこの積分の典型的な
求め方は多変数化して極座標へと変換して
積分を求めることである．
もう一つ筆者が客観的に見て簡単と言える方法に
ベータ関数を使う方法がある．
$\Gamma(1/2)$を求めるとガウス積分の値もわかるのだが，
$\Gamma(1/2)$は$B(1/2, 1/2)$を用いて求めることができ，
この値は初等的な積分になっているのである．
ガンマ関数やベータ関数の
基本的な性質は
\cite{4}を参照のこと．
この本はものすごくいい本なのでとてもおすすめする．
ただし日本語訳は絶版の様子なので，
英語訳版なども参照されたい．日本円で千円前後で
入手できるだろう．

まず，被積分関数が偶関数であることから
\[
\int_{0}^{\infty}\exp(-x^2)\  dx
=\frac{\sqrt{\pi}}{2}
\]
であることを証明できれば良いことに注意しよう．
今から紹介する方法はラプラス \cite{1}を参考にしている．
証明には積分の順序交換を用いる．
\begin{proof}[\textbf{ラプラスの方法によるガウス積分の球値}]
$x\exp(-x^2(y^2+1))$を積分する．
\begin{align*}
	\int_{0}^{\infty}\int_{0}^{\infty}x\exp(-x^2(y^2+1))\ dxdy
	&=
	\int_{0}^{\infty}
	\left[
	-\frac{1}{2}\frac{1}{y^2+1}\exp(-x^2(y^2+1))
	\right]_{x=0}^{x=+\infty}\ dy
	\\
	&=
	\int_{0}^{\infty}
	\frac{1}{2}\frac{1}{y^2+1}\ dy=\frac{\pi}{4}
\end{align*}

積分の順序を交換して
\begin{align*}
	\int_{0}^{\infty}\int_{0}^{\infty}x\exp(-x^2(y^2+1))\ dydx
	&=
	\int_{0}^{\infty}\int_{0}^{\infty}x\exp(-(xy)^2)\exp(-x^2)\ dydx
	\\
	&=
	\int_{0}^{\infty}\exp(-x^2)\int_{0}^{\infty}x\exp(-(xy)^2)\ dydx
\end{align*}
となることがわかる．
さて，
$\int_{0}^{\infty}x\exp(-(xy)^2)\ dy$について$t=xy$と
置換積分をすると
\begin{align*}
	\int_{0}^{\infty}x\exp(-(xy)^2)\ dy
	&=
	\int_{0}^{\infty}\exp(-t^2)\ dt
\end{align*}
なので
\begin{align*}
	\int_{0}^{\infty}\int_{0}^{\infty}x\exp(-x^2(y^2+1))\ dydx
	&=
	\int_{0}^{\infty}\exp(-x^2)\int_{0}^{\infty}\exp(-t^2)\ dtdx
	\\	
	&=
	\int_{0}^{\infty}\exp(-x^2)\ dx\int_{0}^{\infty}\exp(-t^2)\ dt
	\\
	&=
	\left(\int_{0}^{\infty}\exp(-x^2)\ dx\right)^2
\end{align*}

よって
\[
\left(\int_{0}^{\infty}\exp(-x^2)\ dx\right)^2=\frac{\pi}{4}
\]
なので，
\[
\int_{0}^{\infty}\exp(-x^2)\ dx=\frac{\sqrt{\pi}}{2}
\]
がわかる．
\end{proof}

ガウス積分に関連して以下のことがわかる．
\begin{prop}
$a>0$のとき
\[
\int_{-\infty}^{\infty}\exp(-ax^2)\ dx=\sqrt{\frac{\pi}{a}}
\]
\end{prop}

この等式
\[
\int_{-\infty}^{\infty}\exp(-ax^2)\ dx=\sqrt{\frac{\pi}{a}}
\]
において$a$を変数とみて微分をして積分と微分の順序を入れ替えると次のことがわかる．

\begin{cor}
$a>0$のとき
\[
\int_{-\infty}^{\infty}x^2\exp(-ax^2)\ dx=\frac{\sqrt{\pi}}{2a\sqrt{a}}
\]
\end{cor}
特に次の確率論的に関連深い等式が得られる．
\begin{cor}
$m\ge 0$
$\sigma>0$のとき
\[
	\int_{-\infty}^{\infty}
	(x-m)^2
	\frac{1}{\sqrt{2\pi\sigma^2}}
	\exp\left(-\frac{(x-m)^2}{2\sigma^2}\right)\ dx
	=
	\sigma
\]
\end{cor}

\section{フーリエ変換}
このセクションでは
正規分布の確率密度関数のフーリエ変換を具体的に決定する．
つまり正規分布の特性関数を求めるということである．
$f_{m, \sigma}: \rr\to [0, \infty)$を$N(m, \sigma)$
の確率密度関数とする．
つまり
\[
f_{m, \sigma}(x)=
\frac{1}{\sqrt{2\pi\sigma^2}}
\exp\left(-\frac{(x-m)^2}{2\sigma^2}\right)
\]
である．
このとき
\[
F_{m, \sigma}(x)=\int_{\rr}f_{m, \sigma}(y)\exp(\sqrt{-1}xy)\ dy
\]
とおくことにする．

まず最初に次を証明する．
\begin{prop}\label{prop:daiji}
\[
\int_{\rr}\exp(-y^2)\exp(\sqrt{-1}xy)\ dy
=
\sqrt{\pi}
\exp\left(-\frac{x^2}{4}\right)
\]
\end{prop}
\begin{proof}
まず
\[
F(x)=\int_{\rr}\exp(-y^2)\exp(\sqrt{-1}xy)\ dy
\]
とおく．
この$F$を微分する．ルベーグの収束定理から
$F$は$\rr$上の$C^{\infty}$
関数であることに注意しよう．
また同じくルベーグの収束定理から
微分と積分の値が交換できることに注意しよう．
\begin{align*}
	\frac{d}{dx}F(x)
	=
	\sqrt{-1}\int_{\rr}y\exp(-y^2)\exp(\sqrt{-1}xy)\ dy
\end{align*}
である．部分積分を用いると
\begin{align*}
	&\int_{\rr}y\exp(-y^2)\exp(\sqrt{-1}xy)\ dy
	\\
	&=
	\left[-\frac{1}{2}\exp(-y^2)\exp(\sqrt{-1}xy)\right]_{y=-\infty}^{y=\infty}
	+\frac{1}{2}\int_{\rr}\sqrt{-1}x\exp(-y^2)\exp(\sqrt{-1}xy)\ dy
	\\
	&=
	\frac{\sqrt{-1}x}{2}\int_{\rr}\exp(-y^2)\exp(\sqrt{-1}xy)\ dy
	=\frac{\sqrt{-1}x}{2}F(x)
\end{align*}
よって，
\[
\frac{d}{dx}F(x)=-\frac{x}{2}F(x)
\]
がわかる．
この微分方程式を解くと
\[
F(x)=F(0)\exp\left(-\frac{x^2}{4}\right)=\sqrt{\pi}\exp\left(-\frac{x^2}{4}\right)
\]
がわかる．
\end{proof}


\begin{thm}
$m\in [0, \infty)$で
$\sigma\in (0, \infty)$
とする．
このとき
\[
\int_{\rr}f_{m, \sigma}(y)\exp(\sqrt{-1}xy)\ dy
=
\exp\left(\sqrt{-1}mx-\frac{\sigma^2x^2}{2}\right)
\]
\[
F_{m, \sigma}(x)=\exp\left(\sqrt{-1}mx-\frac{\sigma^2x^2}{2}\right)
\]
となる．
特に
\[
F_{0, \sigma}=\exp\left(-\frac{\sigma^2x^2}{2}\right)
\]
であり，
\[
F_{0, 1}=\exp\left(-\frac{x^2}{2}\right)
\]
である．
\end{thm}
\begin{proof}
まず，
\begin{align*}
F_{m, \sigma}(x)
&=
\int_{\rr}f_{m, \sigma}(y)\exp(\sqrt{-1}xy)\ dy\\
&=
\frac{1}{\sqrt{2\pi\sigma^2}}
\int_{\rr}\exp\left(-\frac{(y-m)^2}{2\sigma^2}\right)
\exp(\sqrt{-1}xy)\ dy
\end{align*}
である．
$t=y-m$と
変数変換すると
\begin{align*}
F_{m, \sigma}(x)
&=
\frac{1}{\sqrt{2\pi\sigma^2}}
\int_{\rr}\exp\left(-\frac{t^2}{2\sigma^2}\right)
\exp(\sqrt{-1}x(t+m))\ dt
\\
&=
\exp\left(\sqrt{-1}xm\right)
\frac{1}{\sqrt{2\pi\sigma^2}}
\int_{\rr}\exp\left(-\frac{t^2}{2\sigma^2}\right)
\exp(\sqrt{-1}xt)\ dt\\
&=
\exp\left(\sqrt{-1}xm\right)
F_{0, \sigma}(x)
\end{align*}
なので結局
\[
F_{0, \sigma}(x)=\exp\left(-\frac{\sigma^2x^2}{2}\right)
\]
を示せば良い．
\[
F_{0, \sigma}(x)=
\int_{\rr}\exp\left(-\frac{y^2}{2\sigma^2}\right)
\exp(\sqrt{-1}xy)\ dy
\]
ここで
\[
z=\frac{y}{\sqrt{2\sigma^2}}
\]
と変数変換し，
$F(x)=\sqrt{\pi}\exp(-x^2/4)$とすると
命題 \ref{prop:daiji}から
\begin{align*}
F_{0, \sigma}(x)
&=
\frac{1}{\sqrt{\pi}}
\int_{\rr}\exp\left(-z^2\right)
\exp(\sqrt{-1}(\sqrt{2\sigma^2}x)z)\ dy\\
&=
\frac{1}{\sqrt{\pi}}F(\sqrt{2\sigma^2}x)
=
\exp\left(-\frac{\sigma^2x^2}{2}\right)
\end{align*}
を得るから定理が示された．
\end{proof}








\begin{thebibliography}{9}
\bibitem{1}
P. I. Laplace 著，伊藤清，樋口順四郎約・解説
ラプラス確率論，
現代数学の系譜１２，1986, 共立出版．
\bibitem{2}
リチャード・ファインマン著，
大貫昌子約，
ご冗談でしょう，ファインマンさん I, 
1986, 岩波書店．
\bibitem{3}
F. S. Woods, 
\textit{Advanced calculus}, 
1954, GINN AND COMPANY. 
\bibitem{4}
E. アルティン著, 上野健爾訳・解説, 
ガンマ関数入門, 
2002, 
日本評論社．
\end{thebibliography}

「知識は公共財」


			\end{document}



\begin{comment}
[集合のやつは$\mid$を使う。]
[真なる包含関係]
\subsetneqq
[可換図式]
\[\xymatrix{
&   & (2) \ar@{=>}[rd]&  &  &  & (5) \ar@{=>}[rd]&\\ 
& (1)\ar@{=>}[ru] \ar@{=>}[rd]&  &(4)\ar@{=>}[r]&(7)& (1)\ar@{=>}[ru] \ar@{=>}[rd]& & (7)&\\ 
&   & (3) \ar@{=>}[ur]&  &  &  &(6) \ar@{=>}[ur]&
}\]

[\bigin \endを使わない方法]

{\prop[aa] aaa}
で
\begin{prop}[aa]
aaa
\end{prop}
と同じ\df \thm \proof でも同様

[直和]
$\amalg$

[同値関係でよく使う波線]
\sim

[引数付きの新命令]
\newcommand{\prn}[1]{\|#1\|_p}

[番号のリセット]

\setcounter{equation}{0}


[字体の変更]

{\rm hogehoge}と使う。

\rm ローマン
\it イタリック
\sl 斜体

[supp]

\newcommand{\supp}{\text{{\rm supp}}}

[中央揃えの改行]

	\begin{eqnarray*}
		hoge1&=&hogehoge
		hoge2&=&hogehofe

	\end{eqnarray*}

&&の部分で揃えられる。


[\tag]
\begin{equation}
f(x) = ax^2 + bx + c \tag{1.2.3}
\end{equation}



[場合分け]

	\begin{cases}
		hoge1 & \text{状況１}\\
		hoge2 & \text{状況２}\\
		・
		・
		・
	\end{cases}



\end{comment}





